\section{Related work}
\label{sec:related-work}

\subsection{Sequence-analysis toolkits}

Biopython \citep{cock2009biopython} is the de facto standard Python library
for sequence parsing, translation, and protein-parameter calculation, and is
used directly by a large fraction of bioinformatics scripting in Python.
EMBOSS \citep{rice2000emboss} provides an equivalent, broader suite as
command-line Unix tools rather than a library, including its own
isoelectric-point calculator (\texttt{iep}) and pairwise alignment tools. Both
are mature, well-validated, and neither is superseded by the present work:
\bsi\ reimplements a comparable set of sequence and protein calculations
independently, in pure Python with no required dependencies
(Section~\ref{sec:core-calculations}), and cross-validates that
reimplementation against Biopython directly rather than depending on it, for
reasons discussed in Section~\ref{sec:architecture}. Where the two
deliberately disagree -- most notably on isoelectric point, where no single
pKa scale is universally agreed and independent calculators are documented
to differ by several tenths of a pH unit \citep{kozlowski2016ipc} -- \bsi\
reports the method and conditions used rather than asserting an
unqualified value, and its test suite documents the disagreement rather than
concealing it.

Neither Biopython nor EMBOSS provides an interactive desktop interface, and
neither couples sequence analysis to structure retrieval: a user who wants
both must write connecting code themselves. Several general-purpose
scientific notebook environments (Jupyter-based pipelines built on Biopython)
fill this gap on a per-project basis, but do not constitute a reusable,
tested, packaged tool.

\subsection{Structure databases and prediction resources}

The RCSB Protein Data Bank \citep{berman2000pdb} is the primary repository of
experimentally determined structures and exposes a sequence-search API that
\bsi\ uses directly. The AlphaFold Protein Structure Database
\citep{varadi2022alphafolddb}, built on the AlphaFold prediction method
\citep{jumper2021alphafold}, and the ESM Atlas fold service, built on the
ESMFold language model \citep{lin2023esmfold}, both provide predicted
structures for sequences with no experimental structure, at the cost of a
confidence estimate (predicted local distance difference test, pLDDT) in
place of the crystallographic or cryo-EM validation an experimental
structure carries. \bsi\ does not implement, retrain, or redistribute any of
these prediction methods; it is a client that queries their public APIs,
distinguishes their confidence metric from the superficially similar
crystallographic B-factor reported by experimental structures
(Section~\ref{sec:validation}), and attributes results to the resource that
produced them. None of the three resources' own interfaces performs the
identity check described in Section~\ref{sec:m0m4} against an
externally supplied query sequence; that check is outside their scope as
databases, and is precisely the gap this software is built to close for a
client performing automated or semi-automated retrieval.

\subsection{Structure visualisation and analysis tools}

PyMOL \citep{pymol}, UCSF ChimeraX \citep{pettersen2021chimerax}, and web-based
viewers such as 3Dmol.js \citep{rego20153dmol} provide structure
visualisation and, in the case of PyMOL and ChimeraX, extensive structural
analysis capability once a coordinate file has been obtained. \bsi\ embeds
3Dmol.js for in-application viewing and can hand a retrieved structure to an
external viewer such as PyMOL or ChimeraX when one is installed
(Section~\ref{sec:gui}), rather than duplicating their visualisation or
analysis functionality. Its contribution is upstream of visualisation:
deciding, before a structure is opened in any viewer, whether it is the
structure the user actually asked for.

\subsection{Positioning}

\bsi\ does not claim algorithmic novelty in its sequence alignment, which
uses standard Needleman--Wunsch and semi-global dynamic programming
\citep{needleman1970}, nor in its thermodynamic and physicochemical
calculations, which follow established published methods
(Sections~\ref{sec:core-calculations}--\ref{sec:validation}). Its
contribution is threefold and specifically evidentiary rather than
algorithmic: (i) an explicit, thresholded classification of
sequence-to-structure correspondence (M0--M4) applied uniformly across
three heterogeneous structural resources, in place of the implicit
binary success/failure result common to ad hoc retrieval scripts and to the
initial prototype of this software, whose benchmark reported a hit for every query; (ii) an integrated workflow in which
sequence analysis, structure retrieval, and validation share one
configuration, one cache, and one reproducible project record, rather than
requiring a user to coordinate several independent tools by hand; and (iii)
a body of evidence -- an offline test suite, deterministic fault injection,
cross-validation against Biopython, and a measurement of how often naive
retrieval returns the wrong structure for \N{ds-n} randomly sampled proteins,
judged against independently curated cross-references, together with an
adversarial benchmark -- supporting the reliability of that classification and
reporting its cost and its limits.
